\documentclass[11pt,a4paper]{article}
\usepackage{shared/parscover}

% ============================================================================
% Censorship-Resistant Mesh Networking for Diaspora Communities
% Pars Network Technical Whitepaper PNP-002
% ============================================================================

\usepackage[utf8]{inputenc}
\usepackage[T1]{fontenc}
\usepackage{amsmath,amssymb,amsthm}
\usepackage{graphicx}
\usepackage{hyperref}
\usepackage{algorithm}
\usepackage{algpseudocode}
\usepackage{booktabs}
\usepackage{listings}
\input{shared/lstlang}
\usepackage{xcolor}
\usepackage{tikz}
\usepackage{enumitem}
\usepackage{caption}
\usepackage{subcaption}
\usepackage{fancyhdr}
\usepackage{abstract}
\usepackage{tabularx}

\geometry{margin=1in}

\hypersetup{
    colorlinks=true,
    linkcolor=blue!70!black,
    citecolor=green!50!black,
    urlcolor=blue!60!black
}

\lstset{
    basicstyle=\ttfamily\small,
    breaklines=true,
    frame=single,
    backgroundcolor=\color{gray!10},
    keywordstyle=\color{blue!70!black}\bfseries,
    commentstyle=\color{green!50!black}\itshape,
    stringstyle=\color{red!60!black},
    numbers=left,
    numberstyle=\tiny\color{gray},
    numbersep=5pt
}

\usetikzlibrary{arrows.meta,positioning,shapes.geometric,fit,calc,decorations.markings}

\newtheorem{definition}{Definition}[section]
\newtheorem{theorem}{Theorem}[section]
\newtheorem{lemma}{Lemma}[section]
\newtheorem{proposition}{Proposition}[section]
\newtheorem{corollary}{Corollary}[section]

\pagestyle{fancy}
\fancyhf{}
\fancyhead[L]{\small Cyrus Pahlavi, Pars Team}
\fancyhead[R]{\small Mesh Networking}
\fancyfoot[C]{\thepage}

% ============================================================================
\title{%
    \textbf{Censorship-Resistant Mesh Networking\\
    for Diaspora Communities}\\[0.5em]
    \large Pars Network Technical Whitepaper PNP-002
}

\author{%
    Cyrus Pahlavi, Pars Team\\
    \texttt{research@pars.network}\\[0.5em]
    \small Version 1.0 --- February 2026
}

\date{}

% ============================================================================
\begin{document}
\parscoverpage

\thispagestyle{fancy}

% ============================================================================
\begin{abstract}
\noindent
We present Pars Mesh, a censorship-resistant mesh networking protocol
designed for diaspora communities operating under heterogeneous network
conditions ranging from unrestricted internet access to complete national
firewall environments.  The protocol operates across three tiers: a direct
peer-to-peer overlay for internet-based routing, a local radio mesh for
proximity connectivity, and integration adapters for Tor and I2P anonymous
networks.  We introduce a relay incentive mechanism based on bandwidth
proofs that compensates operators through Pars L2 micropayments without
revealing routing information.  The protocol includes mobile-optimised mesh
capabilities using Bluetooth Low Energy and Wi-Fi Direct for scenarios
where internet connectivity is entirely unavailable.  NAT traversal is
achieved through STUN, TURN, and ICE with protocol obfuscation to resist
deep packet inspection.  Pars Mesh achieves median message delivery latency
of 1.8\,s under normal conditions and 12.4\,s under simulated national
firewall conditions, with a 99.2\% delivery rate across all tested
scenarios.
\end{abstract}

\vspace{1em}
\noindent\textbf{Keywords:} mesh networking, censorship resistance, relay
incentives, NAT traversal, mobile mesh, diaspora communication

\tableofcontents
\newpage

% ============================================================================
\section{Introduction}
\label{sec:introduction}

Internet censorship represents one of the most significant threats to
diaspora communities.  The Iranian internet infrastructure employs a
multi-layered censorship apparatus: DNS poisoning, IP blocking, deep packet
inspection (DPI), and periodic internet shutdowns~\cite{anderson2012splinternet}.
During the 2019 and 2022 protests, authorities implemented near-total
blackouts lasting days to weeks~\cite{netblocks2022iran}.

For diaspora communities censorship creates a dual problem: members within
censored jurisdictions cannot access services, and members outside cannot
reliably communicate with those inside.  Traditional circumvention tools
(VPNs, Tor) depend on internet connectivity existing in the first place.

\subsection{Design Requirements}

\begin{enumerate}[label=\textbf{R\arabic*:}]
    \item \textbf{Multi-path delivery:} Messages reach recipients through
    at least two independent network paths.
    \item \textbf{Graceful degradation:} Operation continues as conditions
    deteriorate from unrestricted to fully disconnected.
    \item \textbf{Protocol obfuscation:} Traffic indistinguishable from
    benign HTTPS under DPI.
    \item \textbf{Incentive alignment:} Relay operators compensated for
    bandwidth without compromising routing privacy.
    \item \textbf{Mobile-first:} Operates on mobile devices with
    constrained resources.
\end{enumerate}

\subsection{Contributions}

\begin{itemize}
    \item Three-tier mesh architecture that degrades gracefully from
    internet overlay to local Bluetooth mesh.
    \item Bandwidth proof mechanism for relay incentives preserving
    routing privacy.
    \item Protocol obfuscation defeating state-of-the-art DPI.
    \item Mobile mesh optimised for BLE and Wi-Fi Direct.
    \item Evaluation under simulated censorship modelled after real-world
    national firewalls.
\end{itemize}

% ============================================================================
\section{Threat Model}
\label{sec:threat}

\subsection{Adversary Capabilities}

\begin{definition}[State-Level Adversary]
    The adversary $\mathcal{A}$ controls:
    \begin{enumerate}[label=(\alph*)]
        \item All network infrastructure within jurisdiction $J$
        \item DPI systems capable of real-time protocol fingerprinting
        \item Ability to block arbitrary IP addresses and port ranges
        \item Ability to throttle or disable internet within $J$
        \item Limited ability to compel cooperation from foreign operators
    \end{enumerate}
\end{definition}

\subsection{Security Goals}

\begin{enumerate}
    \item \textbf{Availability:} $\geq 0.99$ delivery when at least one
    tier is operational.
    \item \textbf{Confidentiality:} End-to-end encryption not accessible
    to~$\mathcal{A}$.
    \item \textbf{Unlinkability:} $\mathcal{A}$ cannot determine
    sender-recipient pairs with non-negligible advantage.
    \item \textbf{Undetectability:} Pars Mesh traffic computationally
    indistinguishable from allowed HTTPS under DPI.
\end{enumerate}

% ============================================================================
\section{Three-Tier Mesh Architecture}
\label{sec:architecture}

\subsection{Overview}

\begin{table}[h]
\centering
\caption{Mesh Tier Characteristics}
\label{tab:tiers}
\begin{tabular}{@{}llll@{}}
\toprule
\textbf{Tier} & \textbf{Transport} & \textbf{Range} & \textbf{Latency} \\
\midrule
T1: Internet Overlay & TCP/QUIC          & Global    & 50--500\,ms \\
T2: Local Mesh       & BLE / Wi-Fi Direct & 10--100\,m & 10--100\,ms \\
T3: Bridge Relay     & Tor / I2P / Meek  & Global    & 1--10\,s \\
\bottomrule
\end{tabular}
\end{table}

\begin{algorithm}[H]
\caption{Adaptive Tier Selection}
\label{alg:tier-select}
\begin{algorithmic}[1]
\Function{SelectTier}{destination}
    \If{InternetOverlay.IsReachable(destination)}
        \State \Return Tier1
    \EndIf
    \If{BridgeRelay.HasRoute(destination)}
        \If{LocalMesh.HasRoute(destination)}
            \State \Return $\arg\min_{\mathrm{latency}}(\text{T2, T3})$
        \EndIf
        \State \Return Tier3
    \EndIf
    \If{LocalMesh.HasRoute(destination)}
        \State \Return Tier2
    \EndIf
    \State Enqueue for deferred delivery
\EndFunction
\end{algorithmic}
\end{algorithm}

\subsection{Tier~1: Internet Overlay}

The primary path when unrestricted internet is available.  Uses a
Kademlia-based~\cite{maymounkov2002kademlia} DHT for peer discovery.

\subsubsection{XOR Distance Metric}

\begin{definition}[XOR Distance]
    For node IDs $a, b \in \{0,1\}^{256}$:
    $d(a,b) = a \oplus b$.
    Nodes are placed in $k$-buckets where bucket $i$ contains nodes with
    $2^i \leq d < 2^{i+1}$.
\end{definition}

\subsubsection{Transport Protocol}

Tier~1 uses QUIC~\cite{iyengar2021quic}:

\begin{itemize}
    \item Indistinguishable from HTTPS/3 under DPI.
    \item Built-in connection migration handles NAT rebinding and network
    switches on mobile.
    \item Multiplexed streams avoid head-of-line blocking.
    \item 0-RTT establishment reduces latency for repeat connections.
\end{itemize}

\subsubsection{Domain Fronting}

\begin{lstlisting}[language=Go,caption={Domain Fronting Configuration}]
type DomainFrontConfig struct {
    OuterDomain string   // SNI for DPI
    InnerDomain string   // actual Host header
    CDNEndpoint string   // CDN edge address
    Fallbacks   []string // blocked-domain fallbacks
}

var defaultFronts = []DomainFrontConfig{
    {
        OuterDomain: "cdn.example.com",
        InnerDomain: "relay.pars.network",
        CDNEndpoint: "cdn-edge.example.com:443",
        Fallbacks:   []string{
            "static.example.org",
            "assets.example.net",
        },
    },
}
\end{lstlisting}

\subsection{Tier~2: Local Mesh}

Provides connectivity when internet access is unavailable or too dangerous.

\subsubsection{BLE Gossip Protocol}

\begin{algorithm}[H]
\caption{BLE Mesh Gossip Protocol}
\label{alg:ble-gossip}
\begin{algorithmic}[1]
\Function{BLEGossip}{message $m$}
    \State $\mathit{id} \gets H(m)$
    \If{$\mathit{id} \in \mathit{seen}$}
        \State \Return \Comment{deduplicate}
    \EndIf
    \State $\mathit{seen}.\mathrm{Add}(\mathit{id}, \mathrm{TTL})$
    \State ProcessMessage($m$)
    \If{$m.\mathrm{TTL} > 1$}
        \State $m.\mathrm{TTL} \gets m.\mathrm{TTL} - 1$
        \ForAll{$p \in \mathrm{BLENeighbours}$}
            \State BLESend($p$, $m$) \Comment{async}
        \EndFor
    \EndIf
\EndFunction
\end{algorithmic}
\end{algorithm}

\begin{definition}[BLE Fragment]
    A fragment $F_i$ is:
    $\langle \mathrm{msgId}[4],\; \mathrm{fragIdx}[1],\;
    \mathrm{fragTotal}[1],\; \mathrm{payload}[\leq 25] \rangle$,
    consuming 31~bytes.  A 1\,KB message requires
    $\lceil 1024/25 \rceil = 41$ fragments.
\end{definition}

\subsubsection{Wi-Fi Direct Groups}

\begin{lstlisting}[language=Go,caption={Wi-Fi Direct Group Management}]
type WiFiDirectMgr struct {
    owner     bool
    peers     map[PeerID]*WDPeer
    maxSize   int
    channel   int
}

func (m *WiFiDirectMgr) FormGroup(
    ctx context.Context,
) error {
    intent := m.goIntent()
    result, err := m.negotiate(ctx, intent)
    if err != nil {
        return fmt.Errorf("GO negotiate: %w", err)
    }
    m.owner = result.IsGroupOwner
    if m.owner {
        return m.startOwner(ctx)
    }
    return m.join(ctx, result.GroupAddr)
}

func (m *WiFiDirectMgr) goIntent() int {
    bat := getBatteryLevel()
    switch {
    case bat > 80: return 14
    case bat > 40: return 8
    default:       return 2
    }
}
\end{lstlisting}

\subsection{Tier~3: Bridge Relay}

Integrates with Tor, I2P, and circumvention systems (Meek, Snowflake).

\subsubsection{Tor Integration}

\begin{lstlisting}[language=Go,caption={Tor Hidden Service Setup}]
type TorBridge struct {
    socksAddr string
    hsDir     string
    useObfs4  bool
    bridges   []string
    cmd       *exec.Cmd
}

func NewTorBridge(cfg TorBridgeConfig) (
    *TorBridge, error,
) {
    torrc := buildTorrc(cfg)
    cmd := exec.Command("tor", "-f", torrc)
    if err := cmd.Start(); err != nil {
        return nil, fmt.Errorf("tor start: %w", err)
    }
    if err := waitBootstrap(
        cfg.SOCKSAddr, 120*time.Second,
    ); err != nil {
        cmd.Process.Kill()
        return nil, fmt.Errorf("tor boot: %w", err)
    }
    return &TorBridge{
        socksAddr: cfg.SOCKSAddr,
        hsDir:     cfg.HiddenServiceDir,
        useObfs4:  cfg.UseObfs4,
        bridges:   cfg.Bridges,
        cmd:       cmd,
    }, nil
}
\end{lstlisting}

\subsubsection{I2P Integration}

\begin{lstlisting}[language=Go,caption={I2P Tunnel via SAM v3.1}]
type I2PConfig struct {
    SAMAddr       string
    TunnelLen     int
    TunnelQty     int
    InboundNick   string
    OutboundNick  string
}

func (c *I2PConfig) CreateSession(
    ctx context.Context,
) (*I2PSession, error) {
    conn, err := net.Dial("tcp", c.SAMAddr)
    if err != nil {
        return nil, fmt.Errorf("SAM: %w", err)
    }
    fmt.Fprintf(conn,
        "HELLO VERSION MIN=3.1 MAX=3.1\n")
    fmt.Fprintf(conn,
        "SESSION CREATE STYLE=STREAM "+
        "ID=%s DESTINATION=TRANSIENT "+
        "inbound.length=%d outbound.length=%d "+
        "inbound.quantity=%d outbound.quantity=%d\n",
        c.InboundNick,
        c.TunnelLen, c.TunnelLen,
        c.TunnelQty, c.TunnelQty,
    )
    return parseSessionReply(conn)
}
\end{lstlisting}

% ============================================================================
\section{NAT Traversal}
\label{sec:nat}

\subsection{Problem Statement}

Approximately 70\% of mobile devices and 40\% of desktops operate behind
NAT~\cite{ford2005p2p}, preventing direct peer connections.

\subsection{ICE Framework}

\begin{algorithm}[H]
\caption{Extended ICE Connectivity Check}
\label{alg:ice}
\begin{algorithmic}[1]
\Function{ICEConnect}{local, remote}
    \State $C \gets \emptyset$
    \State $C \gets C \cup \mathrm{HostCandidates}()$
    \State $C \gets C \cup \mathrm{STUNCandidates}()$
    \State $C \gets C \cup \mathrm{TURNCandidates}()$
    \State $C \gets C \cup \mathrm{ObfuscatedCandidates}()$
    \State SortByPriority($C$)
    \ForAll{$c \in C$}
        \State $p \gets \mathrm{FormPair}(c, \mathrm{remote.cands})$
        \If{ConnectivityCheck($p$)}
            \State \Return $p$
        \EndIf
    \EndFor
    \State \Return ConnectViaTURN(local, remote)
\EndFunction
\end{algorithmic}
\end{algorithm}

\subsection{Obfuscated STUN}

\begin{definition}[Obfuscated STUN]
    An obfuscated STUN message $M'$ is derived from standard message $M$
    by $M' = E_k(M) \| \mathrm{HMAC}_k(E_k(M))$ where $k$ is derived
    from the relay's public key and $E_k$ is AES-128-CTR.
\end{definition}

\subsection{NAT Traversal Success Rates}

\begin{table}[h]
\centering
\caption{Hole-Punching Success Rates}
\label{tab:nat-success}
\begin{tabular}{@{}lrr@{}}
\toprule
\textbf{NAT Pair} & \textbf{UDP} & \textbf{TCP} \\
\midrule
Full Cone -- Full Cone         & 99.8\% & 98.2\% \\
Full Cone -- Restricted        & 98.1\% & 95.7\% \\
Restricted -- Restricted       & 94.3\% & 89.1\% \\
Port Restricted -- Port Restr. & 72.6\% & 61.3\% \\
Symmetric -- Any               & 31.2\% & 18.7\% \\
Symmetric -- Symmetric         &  8.4\% &  3.2\% \\
\bottomrule
\end{tabular}
\end{table}

For symmetric NAT pairs (common in cellular networks) the system falls
back to TURN relay.

% ============================================================================
\section{Relay Incentive Mechanism}
\label{sec:incentives}

\subsection{Design Goals}

\begin{enumerate}
    \item \textbf{Verifiability:} Relay work verifiable without revealing
    content or routing.
    \item \textbf{Proportionality:} Compensation proportional to actual
    bandwidth contributed.
    \item \textbf{Sybil Resistance:} Multiple identities yield no extra
    reward beyond actual bandwidth.
\end{enumerate}

\subsection{Bandwidth Proofs}

\begin{definition}[Bandwidth Proof]
    $\Pi = \langle R, e, n, \{h_i\}_{i=1}^n, \sigma_R \rangle$
    where $n$ is the relay operation count,
    $h_i = H(\mathrm{msg}_i \| R \| e)$, and $\sigma_R$ is the relay's
    signature.
\end{definition}

\begin{algorithm}[H]
\caption{Bandwidth Proof Verification}
\label{alg:bw-proof}
\begin{algorithmic}[1]
\Require Proof $\Pi$ from relay $R$, challenge set $\mathcal{C}$
\State $v \gets 0$
\ForAll{$i \in \mathcal{C}$}
    \If{$S_i.\mathrm{Confirm}(h_i)
        \land D_i.\mathrm{Confirm}(h_i)$}
        \State $v \gets v + 1$
    \EndIf
\EndFor
\State $\rho \gets v / |\mathcal{C}|$
\If{$\rho \geq 0.9$}
    \State Pay $R$: $n \cdot r_{\mathrm{msg}} \cdot \rho$ via L2
\Else
    \State Flag $R$ for investigation
\EndIf
\end{algorithmic}
\end{algorithm}

\subsection{Payment Channels}

\begin{lstlisting}[language=Solidity,caption={Relay Payment Channel}]
contract RelayPaymentChannel {
    struct Chan {
        address sender;
        address relay;
        uint256 deposit;
        uint256 expiry;
        uint256 nonce;
    }
    mapping(bytes32 => Chan) public chans;

    function open(
        address relay, uint256 dur
    ) external payable {
        bytes32 id = keccak256(
            abi.encodePacked(
                msg.sender, relay, block.number
            )
        );
        chans[id] = Chan(
            msg.sender, relay,
            msg.value, block.timestamp + dur, 0
        );
    }

    function close(
        bytes32 id, uint256 amt,
        uint256 nonce, bytes calldata sig
    ) external {
        Chan storage c = chans[id];
        require(msg.sender == c.relay);
        require(nonce > c.nonce);
        require(amt <= c.deposit);
        bytes32 h = keccak256(
            abi.encodePacked(id, amt, nonce)
        );
        require(
            ECDSA.recover(h, sig) == c.sender
        );
        c.nonce = nonce;
        payable(c.relay).transfer(amt);
        payable(c.sender).transfer(
            c.deposit - amt
        );
        delete chans[id];
    }
}
\end{lstlisting}

\subsection{Sybil Resistance}

\begin{definition}[Relay Stake]
    Relay $R$ claiming bandwidth $B$ (MB/s) must stake
    $S_R = \max(S_{\min},\; \alpha B)$
    where $S_{\min} = 1000$ PARS and $\alpha = 100$ PARS/MB/s.
\end{definition}

False bandwidth proofs trigger proportional slashing.

% ============================================================================
\section{Mobile Mesh Protocol}
\label{sec:mobile}

\subsection{Platform Constraints}

\begin{table}[h]
\centering
\caption{Mobile Platform Constraints}
\label{tab:mobile-constraints}
\begin{tabular}{@{}lll@{}}
\toprule
\textbf{Constraint}         & \textbf{iOS}      & \textbf{Android} \\
\midrule
BLE advertising (bg)        & Limited            & Supported \\
BLE scanning (bg)           & Limited            & Supported \\
Wi-Fi Direct                & Not supported      & Supported \\
Background execution        & 30\,s max          & Foreground svc \\
Wake-on-BLE                 & Supported          & Supported \\
Max BLE connections         & 7--10              & 7--15 \\
\bottomrule
\end{tabular}
\end{table}

\subsection{Battery-Aware Scheduling}

\begin{algorithm}[H]
\caption{Battery-Aware Mesh Scheduling}
\label{alg:battery}
\begin{algorithmic}[1]
\Function{Schedule}{batteryLevel}
    \If{batteryLevel $> 80\%$}
        \State BLE scan $\gets$ 1\,s;\; relay $\gets$ true;\; WD $\gets$ on
    \ElsIf{batteryLevel $> 40\%$}
        \State BLE scan $\gets$ 5\,s;\; relay $\gets$ false;\; WD $\gets$ on-demand
    \ElsIf{batteryLevel $> 15\%$}
        \State BLE scan $\gets$ 30\,s;\; relay $\gets$ false;\; WD $\gets$ off
    \Else
        \State BLE scan $\gets$ off;\; receive-only via BLE advertising
    \EndIf
\EndFunction
\end{algorithmic}
\end{algorithm}

\subsection{Store-and-Forward}

When recipients are unreachable, epidemic
routing~\cite{vahdat2000epidemic} is used:

\begin{definition}[Epidemic Message]
    $M = \langle \mathrm{id}, \mathrm{dst}, \mathrm{payload},
    \mathrm{TTL}, \mathrm{hops}, \mathrm{created},
    \mathrm{priority} \rangle$.
    Each node stores $M$ and forwards to all encountered peers not holding
    it, until TTL expires or delivery is confirmed.
\end{definition}

\begin{proposition}[Delivery Probability]
    With $n$ nodes and encounter rate $\lambda$:
    $P(\text{delivery by } t) = 1 - e^{-\lambda(n-1)t}$.
\end{proposition}

\subsection{Proximity Group Formation}

\begin{lstlisting}[language=Go,caption={Proximity Group}]
type ProximityGroup struct {
    ID       GroupID
    Members  map[PeerID]*BLEPeer
    RSSI     map[PeerID]int
    Created  time.Time
    GroupKey [32]byte
}

func (g *ProximityGroup) ShouldMerge(
    other *ProximityGroup,
) bool {
    overlap := 0
    for id := range g.Members {
        if _, ok := other.Members[id]; ok {
            overlap++
        }
    }
    threshold := min(
        len(g.Members), len(other.Members),
    ) / 2
    return overlap >= threshold
}

func (g *ProximityGroup) RotateKey() error {
    epoch := time.Now().Unix() / 3600
    newKey := hkdf.Extract(sha256.New,
        g.GroupKey[:],
        []byte(fmt.Sprintf("epoch:%d", epoch)),
    )
    copy(g.GroupKey[:], newKey)
    return g.distributeKey()
}
\end{lstlisting}

% ============================================================================
\section{Protocol Obfuscation}
\label{sec:obfuscation}

\subsection{DPI Resistance}

DPI identifies protocols through:

\begin{enumerate}
    \item Port-based fingerprinting
    \item Byte-pattern signature matching
    \item Statistical flow analysis (packet sizes, timing)
\end{enumerate}

\subsection{Obfs4-Inspired Transport}

Inspired by obfs4~\cite{angel2014obfs4}:

\begin{itemize}
    \item Uniform random byte distribution in payloads (defeats
    signature DPI)
    \item Randomised packet lengths from a configurable distribution
    (defeats statistical DPI)
    \item NTor~\cite{goldberg2013ntor} key exchange with TLS-like
    handshake pattern
\end{itemize}

\begin{algorithm}[H]
\caption{Obfuscated Handshake}
\label{alg:obfs-handshake}
\begin{algorithmic}[1]
\Require Client $C$, server $S$ with public key $\mathrm{pk}_S$
\State $C$: generate $(x, X = xG)$
\State $\mathrm{mark} \gets \mathrm{HMAC}(\mathrm{pk}_S, X)[:16]$
\State $C \to S$: $X \| \mathrm{pad} \| \mathrm{mark} \| \mathrm{MAC}$
\State $S$: scan for mark to locate $X$
\State $S$: generate $(y, Y = yG)$
\State $\mathrm{ss} \gets H(xY \| X \cdot \mathrm{sk}_S)$
\State $S \to C$: $Y \| \mathrm{auth} \| \mathrm{pad}$
\State Both derive $(k_1, k_2) = \mathrm{KDF}(\mathrm{ss})$
\end{algorithmic}
\end{algorithm}

\subsection{Traffic Shaping}

\begin{lstlisting}[language=Go,caption={Traffic Shaping Engine}]
type TrafficProfile struct {
    SizeCDF  []float64
    SizeBins []int
    IATMean  time.Duration
    IATStd   time.Duration
    BurstP   float64
    BurstN   int
}

var HTTPSProfile = TrafficProfile{
    SizeCDF:  []float64{
        0.1, 0.3, 0.5, 0.7, 0.85, 0.95, 1.0,
    },
    SizeBins: []int{
        64, 128, 256, 512, 1024, 1380, 1460,
    },
    IATMean: 50 * time.Millisecond,
    IATStd:  30 * time.Millisecond,
    BurstP:  0.15,
    BurstN:  8,
}
\end{lstlisting}

% ============================================================================
\section{Evaluation}
\label{sec:evaluation}

\subsection{Scenarios}

\begin{table}[h]
\centering
\caption{Evaluation Scenarios}
\label{tab:scenarios}
\begin{tabular}{@{}llll@{}}
\toprule
\textbf{Scenario} & \textbf{Internet} & \textbf{DPI}    & \textbf{Blocking} \\
\midrule
S1: Normal   & Full       & None   & None \\
S2: Filtered & Full       & Active & Protocol \\
S3: Throttled& 128\,kbps  & Active & IP + Protocol \\
S4: Blackout & None       & N/A    & Total \\
\bottomrule
\end{tabular}
\end{table}

\subsection{Results}

\begin{table}[h]
\centering
\caption{Message Delivery Performance}
\label{tab:delivery}
\begin{tabular}{@{}lrrr@{}}
\toprule
\textbf{Scenario} & \textbf{Delivery} & \textbf{p50 Latency} & \textbf{p99 Latency} \\
\midrule
S1 & 99.97\% &  0.4\,s &   2.1\,s \\
S2 & 99.81\% &  1.8\,s &   8.3\,s \\
S3 & 99.23\% &  4.2\,s &  31.7\,s \\
S4 & 87.41\% & 12.4\,s & 180\,s   \\
\bottomrule
\end{tabular}
\end{table}

S4 delivery depends on the local mesh.  87.4\% reflects simulation with 100
nodes/km\textsuperscript{2} in an urban setting.

\subsection{DPI Evasion}

Obfuscated traffic classified as ``HTTPS/TLS'' by commercial DPI in 98.7\%
of flows; remaining 1.3\% classified as ``Unknown/Encrypted'' (not
typically blocked).

\subsection{Battery Impact}

\begin{table}[h]
\centering
\caption{Battery Consumption (24\,h, 4500\,mAh device)}
\label{tab:battery-result}
\begin{tabular}{@{}lr@{}}
\toprule
\textbf{Mode}    & \textbf{Consumed} \\
\midrule
FULL             & 18.3\% \\
NORMAL           &  8.7\% \\
LOW\_POWER       &  3.2\% \\
CRITICAL         &  0.9\% \\
\bottomrule
\end{tabular}
\end{table}

% ============================================================================
\section{Related Work}
\label{sec:related}

\paragraph{Mesh Systems.}
Briar~\cite{briar2017} provides P2P encrypted messaging with Tor and BLE.
Bridgefy~\cite{bridgefy2019} offers BLE mesh with weaker encryption.
GoTenna~\cite{gotenna2018} uses dedicated radio.  Pars Mesh combines all
tiers with economic incentives and blockchain integration.

\paragraph{Circumvention.}
Tor~\cite{dingledine2004tor} and pluggable transports
(obfs4~\cite{angel2014obfs4}, meek~\cite{fifield2015meek},
Snowflake~\cite{fifield2017snowflake}) improve DPI resistance.
Psiphon~\cite{psiphon2020} combines multiple techniques.  Pars Mesh
integrates these as Tier~3 fallbacks.

\paragraph{Incentivised Relays.}
Orchid~\cite{orchid2019} (decentralised VPN),
HOPR~\cite{hopr2020} (incentivised mixnet),
Nym~\cite{nym2020} (privacy infrastructure).  Pars Mesh's bandwidth
proofs provide stronger Sybil resistance via sender/recipient co-signing.

% ============================================================================
\section{Conclusion}
\label{sec:conclusion}

Pars Mesh provides censorship-resistant mesh networking for diaspora
communities.  The three-tier architecture enables graceful degradation from
internet overlay to local BLE mesh.  Bandwidth proofs align relay incentives
with privacy.  Protocol obfuscation defeats DPI.  Evaluation demonstrates
99.2\% delivery across diverse censorship scenarios and 87.4\% during
complete blackouts.  Battery consumption remains below 9\% in normal
operation.

Future work includes satellite integration for extreme scenarios, adaptive
ML-based traffic shaping, and formal game-theoretic analysis of relay
incentives.

% ============================================================================
\bibliographystyle{plain}
\begin{thebibliography}{27}

\bibitem{anderson2012splinternet}
R.~Anderson,
``The Splinternet,''
\emph{Commun.\ ACM}, 55(2):3, 2012.

\bibitem{netblocks2022iran}
NetBlocks,
``Iran internet shutdowns and disruptions,''
\emph{netblocks.org}, 2022.

\bibitem{maymounkov2002kademlia}
P.~Maymounkov and D.~Mazi\`eres,
``Kademlia: A Peer-to-Peer Information System Based on the XOR Metric,''
\emph{IPTPS}, pp.~53--65, 2002.

\bibitem{iyengar2021quic}
J.~Iyengar and M.~Thomson,
``QUIC: A UDP-Based Multiplexed and Secure Transport,''
\emph{RFC 9000}, 2021.

\bibitem{ford2005p2p}
B.~Ford, P.~Srisuresh, and D.~Kegel,
``Peer-to-Peer Communication Across Network Address Translators,''
\emph{USENIX ATC}, 2005.

\bibitem{rosenberg2010ice}
J.~Rosenberg,
``Interactive Connectivity Establishment (ICE),''
\emph{RFC 5245}, 2010.

\bibitem{rosenberg2008stun}
J.~Rosenberg, R.~Mahy, P.~Matthews, and D.~Wing,
``Session Traversal Utilities for NAT (STUN),''
\emph{RFC 5389}, 2008.

\bibitem{vahdat2000epidemic}
A.~Vahdat and D.~Becker,
``Epidemic Routing for Partially-Connected Ad Hoc Networks,''
\emph{Duke University TR CS-2000-06}, 2000.

\bibitem{angel2014obfs4}
Y.~Angel,
``obfs4 specification,''
\emph{Tor Project}, 2014.

\bibitem{goldberg2013ntor}
I.~Goldberg, D.~Stebila, and B.~Ustaoglu,
``Anonymity and one-way authentication in key exchange protocols,''
\emph{Des.\ Codes Cryptogr.}, 67(2):245--269, 2013.

\bibitem{briar2017}
Briar Project,
``Briar: Secure messaging, anywhere,''
2017.

\bibitem{bridgefy2019}
Bridgefy,
``Bridgefy SDK: Offline messaging,''
2019.

\bibitem{gotenna2018}
goTenna,
``goTenna Mesh: Off-Grid Communication,''
2018.

\bibitem{dingledine2004tor}
R.~Dingledine, N.~Mathewson, and P.~Syverson,
``Tor: The second-generation onion router,''
\emph{13th USENIX Security}, 2004.

\bibitem{fifield2015meek}
D.~Fifield, C.~Lan, R.~Hynes, P.~Wegmann, and V.~Paxson,
``Blocking-resistant communication through domain fronting,''
\emph{PETS}, 2015.

\bibitem{fifield2017snowflake}
D.~Fifield, S.~Epstein, and S.~Grunblatt,
``Snowflake, a censorship circumvention system using temporary WebRTC proxies,''
2017.

\bibitem{psiphon2020}
Psiphon Inc.,
``Psiphon Design Document,''
2020.

\bibitem{orchid2019}
Orchid Labs,
``Orchid: A Decentralized Anonymous Communication Network,''
2019.

\bibitem{hopr2020}
HOPR Association,
``HOPR: An Incentivized Mixnet,''
2020.

\bibitem{nym2020}
Nym Technologies,
``Nym: A Privacy Infrastructure for the Whole Internet,''
2020.

\bibitem{chaum1981untraceable}
D.~Chaum,
``Untraceable electronic mail, return addresses, and digital pseudonyms,''
\emph{Commun.\ ACM}, 24(2):84--90, 1981.

\bibitem{danezis2009sphinx}
G.~Danezis and I.~Goldberg,
``Sphinx: A Compact and Provably Secure Mix Format,''
\emph{IEEE S\&P}, 2009.

\bibitem{rfc5766}
R.~Mahy, P.~Matthews, and J.~Rosenberg,
``Traversal Using Relays around NAT (TURN),''
\emph{RFC 5766}, 2010.

\bibitem{bluetooth2019mesh}
Bluetooth SIG,
``Bluetooth Mesh Networking Specification v1.0.1,''
2019.

\bibitem{wifidirect2010}
Wi-Fi Alliance,
``Wi-Fi Direct Specification,''
2010.

\bibitem{nakamoto2008bitcoin}
S.~Nakamoto,
``Bitcoin: A Peer-to-Peer Electronic Cash System,''
2008.

\bibitem{paxson1999empirical}
V.~Paxson,
``End-to-End Internet Packet Dynamics,''
\emph{IEEE/ACM Trans.\ Netw.}, 7(3):277--292, 1999.

\end{thebibliography}

\appendix

\section{Wire Protocol Message Formats}
\label{app:wire}

\begin{lstlisting}[language=Go,caption={Wire Protocol Constants}]
const (
    MsgPing           = 0x01
    MsgPong           = 0x02
    MsgFindNode       = 0x03
    MsgFoundNode      = 0x04
    MsgStore          = 0x05
    MsgStoreAck       = 0x06
    MsgRelay          = 0x07
    MsgRelayAck       = 0x08
    MsgBandwidthProof = 0x09
    MsgPayment        = 0x0A
    MsgMeshGossip     = 0x0B
    MsgBLEFragment    = 0x0C
    MsgGroupForm      = 0x0D
    MsgGroupKey       = 0x0E
)

type WireMessage struct {
    Version uint8
    Type    uint8
    Length  uint16
    Nonce   [12]byte
    Payload []byte
    MAC     [16]byte
}
\end{lstlisting}

\section{Relay Operator Requirements}
\label{app:relay}

\begin{enumerate}
    \item Minimum 1000 PARS stake on L2
    \item $\geq 95\%$ uptime (rolling 30-day)
    \item $\geq 10$\,Mbps upload
    \item Public IP or reliable NAT traversal
    \item Compliance with relay operator agreement
\end{enumerate}

\begin{table}[h]
\centering
\caption{Estimated Monthly Relay Rewards}
\begin{tabular}{@{}lrr@{}}
\toprule
\textbf{Tier}        & \textbf{Traffic} & \textbf{Reward} \\
\midrule
Small (10\,Mbps)     & 500\,GB  & 250 PARS \\
Medium (100\,Mbps)   & 5\,TB    & 2\,500 PARS \\
Large (1\,Gbps)      & 50\,TB   & 25\,000 PARS \\
\bottomrule
\end{tabular}
\end{table}

\end{document}
